\begin{problem}[第34届物理竞赛复赛][菲涅尔透镜]{79}
    菲涅尔透镜又称同心圆阶梯透镜，它是由很多个同轴环带套在一起构成的，其迎光面是平面，折射面除中心是一个球冠外，其它环带分别是属于不同球面的球台侧面，其纵剖面如右图所示。这样的结构可以避免普通大口径球面透镜既厚又重的缺点。菲涅尔透镜的设计主要是确定每个环带的齿形（即它所属球面的球半径和球心），各环带都是一个独立的（部分）球面透镜，它们的焦距不同，但必须保证具有共同的焦点（即图中 $F$ 点）。已知透镜材料的折射率为 $n$，从透镜中心 $O$（球冠的顶点）到焦点 $F$ 的距离（焦距）为 $f$（平行于光轴的平行光都能经环带折射后会聚到 $F$ 点），相邻环带的间距为 $d$（$d$ 很小，可忽略同一带内的球面像差；$d$ 又不是非常小，可忽略衍射效应）。求
    \begin{subproblem}
        每个环带所属球面的球半径和球心到焦点的距离；
    \end{subproblem}
    \begin{subproblem}
        该透镜的有效半径的最大值和有效环带的条数。
    \end{subproblem}
\end{problem}